% Appendix E: the first-moment lower bound on the pathwidth of sparse random
% graphs, in plain language (owner's request, 2026-10-31). Source:
% reports/random_graph_first_moment.md (Theorems A, A', A'' and the literature
% verdict); Lean: EmptyPair.lean, FirstMoment.lean, FirstMomentAsymptotics.lean.
% The beta*(c) values: reports/scratch/first_moment/beta_star.log (two codes agree).
\section{The pathwidth of sparse random graphs}
\label{app:random}

The equivalences of Section~\ref{sec:equiv} make every statement about the
pathwidth of a graph a statement about the open stacks of an instance. This
appendix proves one such statement for random graphs, in plain language, and
records what of it is checked in Lean. The result is that a sparse random
graph has pathwidth proportional to its number of vertices, with an explicit
constant. The argument is a first moment, an old tool; what is new here is the
exact form of its deterministic step, the constant it yields, and the
machine check. We found no prior statement of the constant.

\subsection{The statement}

$G(n,p)$ is the random graph on $n$ labelled vertices in which each of the
$\binom n2$ possible edges is present with probability $p$, independently. We
take $p=c/n$, so that a vertex has about $c$ neighbours. Write
$\mathrm{pw}$ for pathwidth, $\ln$ for the natural logarithm, and
\[
  \phi(x) = -x\ln x \quad(\phi(0)=0),\qquad
  H(x,y,z) = \phi(x)+\phi(y)+\phi(z)
\]
for the entropy of a probability vector with three entries.

\begin{theorem}[first moment, finite form]
\label{thm:firstmoment}
Let $c>0$ and $0<\beta<1$, and put $a=(1-\beta)/2$ and
\[
  \delta(\beta,c) \;=\; c\,a^{2} \;-\; H(a,a,1-2a).
\]
For every integer $n\ge\max(3,c)$,
\[
  \Pr\bigl[\mathrm{pw}(G(n,c/n))\le\beta n\bigr]
  \;\le\; \exp\bigl(-\delta(\beta,c)\,n+3\ln n+c\,a\bigr)
  \;\le\; e^{c/2}\,n^{3}\,e^{-\delta(\beta,c)\,n}.
\]
\end{theorem}

So whenever $\delta(\beta,c)>0$, the probability that the pathwidth is as small
as $\beta n$ is exponentially small, and the pathwidth is at least $\beta n$
with high probability. The next statement says for which $\beta$ that happens.

\begin{theorem}[the admissible region]
\label{thm:betastar}
If $c\le4\ln2$ there is no $\beta\in(0,1)$ with $\delta(\beta,c)>0$. If
$c>4\ln2$ there is a number $\beta^{*}(c)\in(0,1)$ such that
$\delta(\beta,c)>0$ exactly when $\beta<\beta^{*}(c)$; it is the root in
$(0,1)$ of
\[
  -(1-\beta)\ln\tfrac{1-\beta}{2}-\beta\ln\beta \;=\; \tfrac{c}{4}(1-\beta)^{2}.
\]
Hence $\mathrm{pw}(G(n,c/n))\ge\beta n$ with probability $1-O(n^{3}e^{-\delta n})$
for every $\beta<\beta^{*}(c)$, and as $c\to\infty$,
$1-\beta^{*}(c)=\bigl(4\ln c-4\ln\ln c+O(1)\bigr)/c$.
\end{theorem}

Table~\ref{tab:betastar} gives $\beta^{*}(c)$ at a few values of $c$, beside the
only published explicit lower bound we know for a width parameter of
$G(n,c/n)$, the treewidth bound $1-\sqrt{9\ln3/(2c)}$ of
\citet{PerarnauSerra2014}, which is a lower bound on pathwidth too since
treewidth never exceeds pathwidth.

\begin{table}[htbp]
\centering
\caption{The constant $\beta^{*}(c)$ of Theorem~\ref{thm:betastar}, to four
decimals, and the treewidth bound of \citet{PerarnauSerra2014} at the same
$c$ (blank where that bound is not positive). Two independent programs agree
on every digit of $\beta^{*}$.}
\label{tab:betastar}
\begin{tabular}{@{}lrrrrrrrr@{}}
\toprule
$c$ & 3 & 4 & 6 & 8 & 10 & 20 & 50 & 100 \\
\midrule
$\beta^{*}(c)$ & 0.0087 & 0.0613 & 0.1706 & 0.2628 & 0.3371 & 0.5535 & 0.7606 & 0.8574 \\
$1-\sqrt{9\ln3/(2c)}$ & & & 0.0918 & 0.2136 & 0.2966 & 0.5027 & 0.6855 & 0.7776 \\
\bottomrule
\end{tabular}
\end{table}

The threshold $4\ln2\approx2.77$ in Theorem~\ref{thm:betastar} belongs to the
method, not to the graphs: it is the point at which an edge-free split of the
vertices into two equal halves stops being likely, and below it the first
moment simply says nothing. The pathwidth of $G(n,c/n)$ is in fact linear in
$n$ for every $c>1$ \citep{LeeLeeOum2012}, with a constant that no one has
made explicit.

\subsection{The proof, in four steps}

\paragraph{Step 1: a narrow layout leaves an empty pair.}
Put the vertices in a line and look at any position $i$. The first $i+1$
vertices are the prefix; the rest are the suffix. A suffix vertex is
\emph{active} if it has a neighbour in the prefix. The vertex separation of
the line is the largest number of active vertices at any position, and by
Theorem~\ref{thm:vs} the smallest vertex separation over all lines is the
pathwidth. Now take a line of vertex separation $k$ and any position $i$. The
suffix has $n-i-1$ vertices, at most $k$ of them active, so at least $n-i-1-k$
suffix vertices have no neighbour at all in the prefix. Call that set $R$ and
the prefix $A$. Then $A$ and $R$ are disjoint, $|A|=i+1$, $|R|\ge n-i-1-k$, and
there is no edge between them: an \emph{empty pair}. This holds at every
position of every line, so a graph of pathwidth $k$ has an empty pair with
$|A|=t$ and $|R|\ge n-t-k$ for every $t\le n$. With $t=\lceil(n-k)/2\rceil$
both parts have about $(n-k)/2$ vertices~\leanref{emptypair}.

\paragraph{Step 2: count the graphs that have one.}
Fix a pair of disjoint sets $A$ and $R$ with $|A|=t$ and $|R|=t-1$. In
$G(n,p)$ the $t(t-1)$ possible edges between them are all absent with
probability $(1-p)^{t(t-1)}$. There are $\binom nt\binom{n-t}{t-1}$ such
pairs. A graph of pathwidth at most $k$ has at least one of them empty, so by
the union bound
\[
  \Pr[\mathrm{pw}\le k]\;\le\;\binom nt\binom{n-t}{t-1}\,(1-p)^{t(t-1)},
  \qquad t=\lceil(n-k)/2\rceil .
\]
This is the whole probabilistic content of the proof, and it is a Lean
theorem over the finite sum of the weights $p^{|E|}(1-p)^{\binom n2-|E|}$ of
all graphs on $n$ vertices, with no probability model needed; the count of
graphs with a given empty pair is exact there, $2^{\binom n2-t(t-1)}$ of
them~\leanref{firstmoment}.

\paragraph{Step 3: estimate the count.}
Two standard inequalities turn the right-hand side into an exponential in
$n$. The number of ways to choose the pair is the trinomial coefficient
$n!/\bigl(t!\,(t-1)!\,(n-2t+1)!\bigr)$, and a trinomial coefficient is at most
$\exp(nH)$ with $H$ the entropy of the three proportions $t/n$, $(t-1)/n$ and
$(n-2t+1)/n$ \citep[Theorem~11.1.3]{CoverThomas2006}: expand
$(x+y+z)^{n}=1$ at $x=t/n$, $y=(t-1)/n$, $z=(n-2t+1)/n$ and keep one term of
the sum. And $(1-p)^{m}\le e^{-pm}$ because $1-p\le e^{-p}$. With $p=c/n$ the
second factor is at most $\exp(-c\,t(t-1)/n)$.

\paragraph{Step 4: the bookkeeping.}
The three proportions are not exactly $a$, $a$ and $1-2a$: with
$k=\lfloor\beta n\rfloor$, each is within $1/n$ of its target. A shift of $d$
in one argument of $\phi$ changes $\phi$ by at most $\phi(d)$ when $d\le1/e$,
and $\phi(1/n)=(\ln n)/n$; three shifts cost $3\ln n$ in the exponent, which
is the $3\ln n$ of the theorem and the reason for $n\ge3$. Likewise
$t(t-1)\ge a^{2}n^{2}-an$, and the $-an$ costs the term $c\,a$. Multiplying
the two estimates of Step~3 into the bound of Step~2 gives
$\exp(nH(a,a,1-2a)-ca^{2}n+3\ln n+ca)$, which is Theorem~\ref{thm:firstmoment};
$a\le1/2$ gives its second form. Nothing is hidden in an error term: the
inequalities of Step~3 are exact bounds on the two factors, and the two
corrections are the ones just named~\leanref{firstmomentA}.

\paragraph{The admissible region.}
On $(0,\tfrac12]$ the function $g(a)=H(a,a,1-2a)-ca^{2}$ is strictly concave,
tends to $0$ as $a\to0$, and equals $\ln2-c/4$ at $a=\tfrac12$. If
$c\le4\ln2$ it is nonnegative throughout, so $\delta=-g$ is never positive.
If $c>4\ln2$ it has exactly one zero $a^{*}$ in $(0,\tfrac12)$, is positive
before it and negative after it, so $\delta(\beta,c)>0$ exactly when
$a=(1-\beta)/2>a^{*}$, that is when $\beta<1-2a^{*}=:\beta^{*}(c)$.
Substituting $a=(1-\beta)/2$ into $g(a)=0$ gives the equation of
Theorem~\ref{thm:betastar}, and the asymptotic form follows from the
two-sided estimate $2\ln(1/a)-ca\le g(a)/a\le2\ln(1/a)+2-ca$. This part is
proved on paper only.

\subsection{What is checked, and what is known}

Steps 1, 2, 3 and 4 and the two inequalities of Theorem~\ref{thm:firstmoment}
are Lean theorems with the statement exactly as printed, including the
constant~3; so is the consequence that the probability tends to zero whenever
$\delta(\beta,c)>0$. The entropy bound on trinomial coefficients is not in
Mathlib and is proved there from the binomial theorem. Theorem~\ref{thm:betastar}
is proved on paper, and the values in Table~\ref{tab:betastar} were computed
by two programs written independently.

The argument is published twice in cruder form. \citet{Turner1986} used the
same empty pair, with the count $(en/k)^{2k}$ in place of the entropy, to
prove that almost every $G(n,p)$ with $p$ fixed has bandwidth at least
$n-4\log_{1/(1-p)}n$; since bandwidth is at least pathwidth, our dense
statement implies his and not conversely. \citet{DiazPetitSernaTrevisan2001}
used it for vertex separation on mixing graphs, with a $3^{n}$ count and
Chernoff slack at a single position, which gives no positive constant below
$c\approx227$. Neither states the entropy form, a value of $\beta^{*}(c)$, the
threshold $4\ln2$ or the asymptotics, and we found no later work that does.

One computation says how far the finite inequality is from the truth. On
133 random graphs with 60 vertices and $c=10$ whose pathwidth we computed
exactly, the smallest pathwidth is 25, while Theorem~\ref{thm:firstmoment}'s
first inequality, evaluated exactly rather than through Step~3, puts the
probability of pathwidth at most 23 below $0.35$ and of at most 21 below
$10^{-3}$. The gap between $\beta^{*}(c)$ and the slopes fitted to such data
is of the order of $0.1$ and does not close with $n$; closing it would need a
second moment or a count over structured pairs, not a sharper estimate of
this one.
